\section{Methods}
\label{sec:methods}

\subsection{Notation and problem setup}

The docstring--test--code triangle comprises three artefacts associated with a single Python function: the code $C$, its docstring $D$, and its reference test suite $T$. A closure operator traverses two edges of the triangle, generating a reconstructed artefact from the third; the pipeline is \emph{closed} on this traversal when the reconstruction is semantically equivalent to the original. Table~\ref{tab:notation} summarises the notation used throughout this chapter.

\begin{table}[h]
\centering
\small
\caption{Notation used in the round-trip closure formalism.}
\label{tab:notation}
\begin{tabular}{ll}
\toprule
Symbol & Meaning \\
\midrule
$C, D, T$ & original code, docstring, and reference test suite \\
$C', D', T'$ & pipeline-reconstructed artefacts \\
$L_{\textsf{spec}}, L_{\textsf{test}}, L_{\textsf{code}}$ & Small Language Models assigned to the three stages \\
$L_J$ & external judge SLM (DeepSeek-R1, family-disjoint from pipeline) \\
$\mathcal{M} = \{m_1, \ldots, m_6\}$ & the six pipeline SLMs (see Table~\ref{tab:model_lineup}) \\
$\mathcal{S} = \{\textsf{spec}, \textsf{test}, \textsf{code}\}$ & pipeline stages \\
$a : \mathcal{S} \to \mathcal{M} \cup \{\bot\}$ & cell assignment (mapping stages to models or ablation) \\
$K(T', C)$ & mutation kill rate of test suite $T'$ against code $C$ \\
$R(C', T)$ & reference test pass rate of code $C'$ against tests $T$ \\
$B(D_1, D_2)$ & rescaled BERTScore-F1 between two docstrings \\
$J(x, y) \in \{-1, 0, 1, 2, 3, 4\}$ & judge equivalence rating; $-1$ denotes parse failure \\
$\tau$ & metric-validity threshold (path-specific; $\tau = 0$ in our defaults) \\
$\rho$ & judge-validity threshold ($\rho = 3$: ``equivalent or better'') \\
$\mathds{1}[\cdot]$ & indicator function \\
$N$ & per-cell sample size ($150$ for mono, $75$ for hetero and null) \\
\bottomrule
\end{tabular}
\end{table}

\subsection{Closure paths}

We define three closure paths, each traversing two edges of the docstring--test--code triangle. Path 1 (Code $\to$ Docstring $\to$ Tests) tests whether tests derived from an LLM-generated docstring can still detect defects in the original code. Path 2 (Docstring $\to$ Tests $\to$ Code) tests whether code reconstructed from the original docstring and LLM-generated tests can satisfy the reference test suite. Path 3 (Code $\to$ Tests $\to$ Docstring) tests whether the docstring recovered from tests-generated-from-code still describes the original function. Formally,

\begin{align}
  \text{Path 1 (}C \to D \to T\text{):} \quad
  T' &\;=\; L_{\textsf{test}}\bigl(L_{\textsf{spec}}(C)\bigr), \label{eq:path1}\\
  \text{Path 2 (}D \to T \to C\text{):} \quad
  C' &\;=\; L_{\textsf{code}}\bigl(D,\; L_{\textsf{test}}(D)\bigr), \label{eq:path2}\\
  \text{Path 3 (}C \to T \to D\text{):} \quad
  D' &\;=\; L_{\textsf{spec}}\bigl(L_{\textsf{test}}(C)\bigr). \label{eq:path3}
\end{align}

Path 2 explicitly conditions $L_{\textsf{code}}$ on both the original docstring $D$ and the freshly-generated tests $T'$ --- this matches the strongest reconstruction context available to the pipeline. Algorithm~\ref{alg:path_execution} presents the unified execution procedure.

\begin{algorithm}[h]
\caption{Round-trip closure path execution}
\label{alg:path_execution}
\begin{algorithmic}[1]
\Require cell $a = (L_{\textsf{spec}}, L_{\textsf{test}}, L_{\textsf{code}})$, sample $s = (C, D, T)$
\Ensure closure result $r_p$ for each path $p \in \{1, 2, 3\}$
\For{$p \in \{1, 2, 3\}$}
  \If{$p = 1$} \Comment{$C \to D \to T$}
    \State $D' \leftarrow L_{\textsf{spec}}(C)$
    \State $T' \leftarrow L_{\textsf{test}}(D', \textsf{entry\_point}, \textsf{signature})$
    \State $T'_{\textsf{filt}} \leftarrow \textsf{TestFilter}(T', C)$ \Comment{Algorithm~\ref{alg:test_filter}}
    \State $K \leftarrow \textsf{MutationKillRate}(T'_{\textsf{filt}}, C)$
    \State $J \leftarrow \textsf{Judge}(D, D', \text{kind}=\textsc{docstring})$
    \State $r_1 \leftarrow (\text{metric}=K,\; \text{judge}=J)$
  \ElsIf{$p = 2$} \Comment{$D \to T \to C$}
    \State $T' \leftarrow L_{\textsf{test}}(D, \textsf{entry\_point}, \textsf{signature})$
    \State $C' \leftarrow L_{\textsf{code}}(D, T', \textsf{entry\_point}, \textsf{signature})$
    \State $R \leftarrow \textsf{ReferencePassRate}(T, C')$
    \State $J \leftarrow \textsf{Judge}(C, C', \text{kind}=\textsc{code})$
    \State $r_2 \leftarrow (\text{metric}=R,\; \text{judge}=J)$
  \Else \Comment{$p = 3$, $C \to T \to D$}
    \State $T' \leftarrow L_{\textsf{test}}(C)$
    \State $D' \leftarrow L_{\textsf{spec}}(T')$
    \State $B \leftarrow \textsf{BERTScore}(D, D')$
    \State $J \leftarrow \textsf{Judge}(D, D', \text{kind}=\textsc{docstring})$
    \State $r_3 \leftarrow (\text{metric}=B,\; \text{judge}=J)$
  \EndIf
\EndFor
\State \Return $[r_1, r_2, r_3]$
\end{algorithmic}
\end{algorithm}

\subsection{Closure validity decision}

A round-trip closure is considered \emph{valid} when both the automated metric and the judge SLM agree on the artefact equivalence. Formally, for Path 1 we require

\begin{equation}
  \text{valid}_1(C, D, D', T')
  \;=\;
  \mathds{1}\!\bigl[\, K(T', C) > \tau \,\bigr]
  \;\wedge\;
  \mathds{1}\!\bigl[\, J(D, D') \geq \rho \,\bigr],
  \label{eq:valid1}
\end{equation}

and analogously for Paths 2 and 3:

\begin{equation}
  \text{valid}_2 = \mathds{1}[R(C', T) > 0] \wedge \mathds{1}[J(C, C') \geq \rho],
  \label{eq:valid2}
\end{equation}
\begin{equation}
  \text{valid}_3 = \mathds{1}[B(D, D') > 0] \wedge \mathds{1}[J(D, D') \geq \rho].
  \label{eq:valid3}
\end{equation}

The two-signal \emph{strict-AND} policy is deliberate: it prevents metric noise or judge parse failures from spuriously inflating closure rates. We adopt the strict policy because the alternative --- single-signal validity --- was the failure mode responsible for inflated closure rates observed in the 30-function pilot; the strict policy also makes the false-closure rate (Section~5) well-defined.

Every closure row is additionally categorised by Algorithm~\ref{alg:validity} into one of five decision-reason labels that exhaust the two-signal decision space. The disagreement categories (\textsf{false\_closure\_candidate} and \textsf{metric\_false\_negative}) do not count as valid closures, but they become the analytical hooks in Section~5 for characterising \emph{where} the automated metric and the judge diverge.

\begin{algorithm}[h]
\caption{Closure validity decision (implements Eq.~\ref{eq:valid1}--\ref{eq:valid3})}
\label{alg:validity}
\begin{algorithmic}[1]
\Require metric value $m$, judge rating $J \in \{-1, 0, 1, 2, 3, 4\}$, thresholds $\tau$, $\rho$
\Ensure validity $\in \{\text{valid}, \text{invalid}\}$, \textsf{decision\_reason} $\in$ five categories
\If{$m$ is missing \textbf{or} $J = -1$}
  \State \Return $(\text{invalid},\; \textsf{structural\_NA})$ \Comment{Ablation cells or empty artefacts}
\EndIf
\If{$m > \tau$ \textbf{and} $J \geq \rho$}
  \State \Return $(\text{valid},\; \textsf{both\_agree\_valid})$
\ElsIf{$m \leq \tau$ \textbf{and} $J < \rho$}
  \State \Return $(\text{invalid},\; \textsf{both\_agree\_invalid})$
\ElsIf{$m > \tau$ \textbf{and} $J < \rho$}
  \State \Return $(\text{invalid},\; \textsf{false\_closure\_candidate})$ \Comment{Metric over-credits}
\Else \Comment{$m \leq \tau$ and $J \geq \rho$}
  \State \Return $(\text{invalid},\; \textsf{metric\_false\_negative})$ \Comment{Metric under-credits}
\EndIf
\end{algorithmic}
\end{algorithm}

Algorithm~\ref{alg:validity} is implemented as the pure function \texttt{closure\_decision.decide\_validity} in the replication package, with 43 unit tests covering all five decision reasons, the structural-NA edge cases (NaN metric, judge $= -1$, \texttt{None}), and the boundary values of $\tau$ and $\rho$.

\subsection{Test-filter validity gate}

Before computing the mutation kill rate, we discard generated tests that fail on the original code $C$. This avoids penalising mutants for surviving tests that were spuriously wrong to begin with. Formally, given a candidate test set $T' = \{t_1, t_2, \ldots\}$, the filtered suite is

\begin{equation}
  T'_{\textsf{filt}}(C)
  \;=\;
  \bigl\{\, t \in T' \;:\; \textsf{exec}(t, C) = \textsf{pass} \,\bigr\}.
  \label{eq:filter}
\end{equation}

The filter is implemented as Algorithm~\ref{alg:test_filter}, which additionally emits a machine-readable \textsf{filter\_reason} label recording \emph{why} the filter output was empty when it was. This label appears in every closure row in the results TSV and is used in Section~5 to distinguish rows where the LLM produced no tests at all (\textsf{empty\_input}) from rows where every generated test failed on the ground truth (\textsf{all\_dropped}) --- a distinction that matters for the phi-4-in-test-stage analysis in Section~5.3.

\begin{algorithm}[h]
\caption{Test-filter validity gate (implements Eq.~\ref{eq:filter})}
\label{alg:test_filter}
\begin{algorithmic}[1]
\Require candidate test suite $T'$, original code $C$
\Ensure filtered test suite $T'_{\textsf{filt}} \subseteq T'$, \textsf{filter\_reason} label
\If{$T'$ or $C$ is empty}
  \State \Return $(\emptyset,\; \textsf{empty\_input})$
\EndIf
\State $N \leftarrow$ number of \texttt{def test\_*} functions in $T'$
\If{$N = 0$}
  \State \Return $(\emptyset,\; \textsf{no\_test\_functions})$
\EndIf
\State $T'_{\textsf{filt}} \leftarrow \emptyset$
\For{each test $t$ in $\textsf{split\_test\_functions}(T')$}
  \State $\textsf{status} \leftarrow \textsf{pytest\_subprocess}(t, C, \textsf{timeout} = 10\text{s})$
  \If{$\textsf{status} = \textsf{pass}$}
    \State $T'_{\textsf{filt}} \leftarrow T'_{\textsf{filt}} \cup \{t\}$
  \EndIf
\EndFor
\If{$|T'_{\textsf{filt}}| = 0$}
  \State \Return $(\emptyset,\; \textsf{all\_dropped})$ \Comment{Row will be written with NaN metric}
\EndIf
\State \Return $(T'_{\textsf{filt}},\; \textsf{kept}\_K\_\textsf{of}\_N)$
\end{algorithmic}
\end{algorithm}

\subsection{Closure metrics}

\subsubsection{Path 1: mutation kill rate}

Given the filtered test suite $T'_{\textsf{filt}}$ and the set of mutants $\mathcal{M}_C = \{m_1, \ldots, m_n\}$ generated from $C$, the mutation kill rate is

\begin{equation}
  K(T', C)
  \;=\;
  \frac{\bigl|\{\, m \in \mathcal{M}\emph{C \;:\; \exists t \in T'}{\textsf{filt}}(C),\; \textsf{exec}(t, m) = \textsf{fail} \,\}\bigr|}{\bigl|\mathcal{M}_C\bigr| \;-\; \bigl|\mathcal{M}_C^{\textsf{equiv}}\bigr|}.
  \label{eq:killrate}
\end{equation}

A mutant is \emph{killed} if at least one filtered test fails on it. Equivalent mutants $\mathcal{M}_C^{\textsf{equiv}}$ (mutants whose observable behaviour matches the original) are excluded from the denominator following Chapter~2's protocol. Five mutation operators are applied: arithmetic (e.g. $\texttt{+} \to \texttt{-}$), boundary ($\texttt{<} \to \texttt{<=}$), comparison ($\texttt{<} \to \texttt{>}$), negate\_bool, and return\_none. Up to 15 mutants are generated per function.

\subsubsection{Path 2: reference test pass rate}

Let $T = \{t_1, \ldots, t_k\}$ be the ground-truth reference suite. The pass rate is

\begin{equation}
  R(C', T)
  \;=\;
  \frac{1}{|T|}\sum_{i=1}^{|T|} \mathds{1}\!\bigl[\,\textsf{exec}(t_i, C') = \textsf{pass}\,\bigr].
  \label{eq:passrate}
\end{equation}

In this experiment we use a binary indicator ($R \in \{0, 1\}$: the suite either fully passes or does not) for the Section~4 statistical analysis.

\subsubsection{Path 3: BERTScore F1}

For two docstrings $D_1$ and $D_2$, let $\phi(D)$ denote their RoBERTa-large contextual embeddings. Precision and recall under cosine similarity are

\begin{align}
  P(D_1, D_2) &\;=\; \tfrac{1}{|\phi(D_2)|} \textstyle\sum_{x \in \phi(D_2)} \max_{y \in \phi(D_1)} \cos(x, y), \label{eq:bert_p}\\
  R(D_1, D_2) &\;=\; \tfrac{1}{|\phi(D_1)|} \textstyle\sum_{y \in \phi(D_1)} \max_{x \in \phi(D_2)} \cos(x, y), \label{eq:bert_r}
\end{align}

and the BERTScore F1 we report is

\begin{equation}
  B(D_1, D_2)
  \;=\;
  \frac{2 \cdot P(D_1, D_2) \cdot R(D_1, D_2)}{P(D_1, D_2) + R(D_1, D_2)}.
  \label{eq:bertscore}
\end{equation}

We use baseline-rescaled scores so $B \in (-\infty, 1]$ with $B = 0$ denoting ``no better than random pairing''. In our data, $B$ values on real closure pairs concentrate in $[-0.2, 0.7]$.

\subsubsection{Judge equivalence rating}

The judge SLM $L_J$ = DeepSeek-R1:14B is prompted with a 5-level behaviourally-anchored rubric: 4 = identical, 3 = equivalent (same observable behaviour), 2 = approximately equivalent, 1 = clearly different, 0 = unrelated. A rating of $-1$ denotes parse failure --- the judge is treated as absent for that row. DeepSeek-R1 is deliberately chosen because it belongs to a family (DeepSeek) disjoint from every pipeline SLM, avoiding self-evaluation bias.

\subsection{Design of experiments}

A cell is a stage-to-model assignment $a : \mathcal{S} \to \mathcal{M} \cup \{\bot\}$, where $\bot$ denotes an ablated stage. The DOE consists of 20 pre-registered cells across three strata:

\begin{equation}
  \mathcal{C}_{\textsf{DOE}} \;=\;
  \underbrace{\bigl\{a : a(\textsf{spec}) = a(\textsf{test}) = a(\textsf{code})\bigr\}}_{\text{Mono stratum (6 cells)}}
  \;\cup\;
  \underbrace{\bigl\{a : |\{a(s) : s \in \mathcal{S}\}| = 3\bigr\}}_{\text{Hetero stratum (11 cells)}}
  \;\cup\;
  \underbrace{\mathcal{C}\emph{{\textsf{null}}}}{\text{Null stratum (3 cells)}}.
  \label{eq:doe}
\end{equation}

\begin{table}[t]
\centering
\small
\caption{Pre-registered 20-cell DOE over six open-weight SLMs and three pipeline stages ($6^3 = 216$ full design space, sampled purposively, not by fractional factorial). Three strata: 6 mono, 11 hetero, 3 null. Hypothesis is the rationale recorded at design time.}
\label{tab:doe_summary}
\begin{tabular}{@{}llllll@{}}
\toprule
Cell & Stratum & $L_{\textsf{spec}}$ & $L_{\textsf{test}}$ & $L_{\textsf{code}}$ & Hypothesis \\
\midrule
M1 & mono & llama3.2 & llama3.2 & llama3.2 & Small-dense self-consistency floor (3\,B Meta). \\
M2 & mono & phi4 & phi4 & phi4 & Mid-dense reasoning-tuned mono baseline (14\,B Microsoft). \\
M3 & mono & qwen3.6 & qwen3.6 & qwen3.6 & Dense mono baseline (27\,B Alibaba). \\
M4 & mono & gemma4 & gemma4 & gemma4 & MoE mono baseline (26\,B Google). \\
M5 & mono & mistral-s.\ 3.2 & mistral-s.\ 3.2 & mistral-s.\ 3.2 & Function-calling-tuned dense mono (24\,B Mistral). \\
M6 & mono & qwen3-coder & qwen3-coder & qwen3-coder & Code-specialised MoE self-consistency (30\,B Alibaba). \\
H1 & hetero & phi4 & qwen3-coder & qwen3-coder & Predicate-reasoner at spec plus coder at test/code. \\
H2 & hetero & qwen3.6 & phi4 & qwen3-coder & Dense spec, reasoner tests, coder synth. \\
H3 & hetero & gemma4 & mistral-s.\ 3.2 & qwen3-coder & Cross-family triple: Google $\to$ Mistral $\to$ Alibaba-coder. \\
H4 & hetero & qwen3-coder & qwen3-coder & llama3.2 & Cheap drafter at synth: does a 3\,B model suffice? \\
H5 & hetero & llama3.2 & qwen3-coder & qwen3-coder & Cheap drafter at spec: can a 3\,B docstring survive downstream? \\
H6 & hetero & phi4 & qwen3.6 & phi4 & Phi sandwich with a different-family middle stage. \\
H7 & hetero & qwen3-coder & qwen3.6 & phi4 & Reverse-capability gradient: strongest first, weakest last. \\
H8 & hetero & gemma4 & qwen3-coder & mistral-s.\ 3.2 & MoE $\to$ MoE $\to$ dense: does architecture matching matter? \\
H9 & hetero & qwen3-coder & qwen3.6 & qwen3-coder & Strong-sandwich: MoE bookends dense for synthesis stability. \\
H10 & hetero & mistral-s.\ 3.2 & phi4 & qwen3-coder & Best-of-each-family-stage: Mistral spec + Phi test + Qwen-coder synth. \\
H11 & hetero & phi4 & gemma4 & qwen3-coder & Phi spec + Gemma test + Qwen-coder synth (alt family triple). \\
N1 & null & llama3.2 & llama3.2 & llama3.2 & Word-shuffled input: expected metric $\approx 0$ if metric is sound. \\
N2 & null & (skip) & qwen3-coder & qwen3-coder & Spec-stage ablation: quantifies $L_{\textsf{spec}}$ contribution. \\
N3 & null & qwen3-coder & (skip) & qwen3-coder & Test-stage ablation: quantifies $L_{\textsf{test}}$ contribution. \\
\bottomrule
\end{tabular}
\end{table}


\begin{table}[t]
\centering
\caption{Small Language Models in the Chapter-3 pipeline. All models are open-weight and served via Ollama; total parameter count is $<30$ B per the §1.1 SLM definition.}
\label{tab:model_lineup}
\begin{tabular}{@{}lccccc@{}}
\toprule
Role & Ollama tag & Family & Size (B) & Architecture & Year \\
\midrule
Pipeline & llama3.2:3b & Meta & 3.0 & dense & 2024 \\
Pipeline & phi4:14b & Microsoft & 14.0 & dense & 2024 \\
Pipeline & qwen3.6:27b & Alibaba & 27.0 & dense & 2026 \\
Pipeline & gemma4:26b & Google & 26.0 & MoE & 2026 \\
Pipeline & mistral-small3.2:24b & Mistral & 24.0 & dense & 2025 \\
Pipeline & qwen3-coder:30b & Alibaba-coder & 30.0 & MoE & 2025 \\
Judge & deepseek-r1:14b & DeepSeek & 14.0 & dense & 2025 \\
\bottomrule
\end{tabular}
\end{table}


Table~\ref{tab:doe_summary} presents the complete DOE with the pre-registered hypothesis for each cell. The full 5-model $\times$ 3-stage design space contains $6^3 = 216$ cells; we sample 20 as a hypothesis-driven fractional factorial. The mono stratum establishes each SLM's self-consistency floor as a baseline. The hetero stratum assigns three different SLMs across the three stages based on per-stage strengths derived from Chapter~2's per-operator capability map (phi-4 dominates predicate reasoning; qwen-coder dominates comparison operators). The null stratum comprises three control cells: N1 (word-shuffled input; expected metric $\approx 0$ if the closure metric is sound), N2 (spec-stage ablated; $L_{\textsf{spec}} = \bot$), and N3 (test-stage ablated; $L_{\textsf{test}} = \bot$).

The six pipeline SLMs (Table~\ref{tab:model_lineup}) are all open-weight and under 30~B parameters, ensuring the entire experiment fits within a single Colab~Pro+ A100/T4 monthly allowance. Five distinct model families are represented (Meta, Microsoft, Alibaba dense, Google, Mistral, Alibaba coder). DeepSeek-R1:14B serves as the family-disjoint judge $L_J$.

\subsection{Datasets}

The primary dataset (\emph{core}) is a stratified 150-function sample from HumanEval (45~functions) and MBPP (105~functions), drawn with seed~42. This preserves the MBPP/HumanEval ratio of the 300-function Chapter~2 sample and enables direct cross-paper comparisons via shared per-function difficulty draws.

Two held-out subsets were prepared for contamination sensitivity analysis: a 25-problem LiveCodeBench subset with publication date $\geq$ 2024-12-01 (post-training-cutoff for all evaluated SLMs) and a 50-problem HumanEval-Mutated subset obtained by applying function-rename, docstring-paraphrase, and parameter-permutation transformations to a subset of HumanEval. The held-out subsets were prepared but not fully swept in this study; contamination sensitivity results on them will be reported in an addendum.

\subsection{Statistical methodology}

\subsubsection{Analysis of variance}

For each closure metric $Y$ (kill rate, pass rate, BERTScore stacked across paths), we fit a Type-III analysis-of-variance model

\begin{equation}
  Y_{ij} \;=\; \mu \;+\; \alpha_i \;+\; \beta_j \;+\; \varepsilon_{ij},
  \quad \varepsilon_{ij} \sim \mathcal{N}(0, \sigma^2),
  \label{eq:anova}
\end{equation}

where $\alpha_i$ is the fixed effect of cell $i \in \mathcal{C}_{\textsf{DOE}}$, $\beta_j$ is the fixed effect of sample $j$, and $\varepsilon_{ij}$ is residual error. To further test whether the three closure paths carry non-redundant information about cell performance, we extend Eq.~\ref{eq:anova} to include a path factor and its interaction with cell:

\begin{equation}
  Y_{ijk} \;=\; \mu \;+\; \alpha_i \;+\; \eta_k \;+\; (\alpha \cdot \eta)_{ik} \;+\; \beta_j \;+\; \varepsilon_{ijk},
  \label{eq:anova_interact}
\end{equation}

where $\eta_k$ is the fixed effect of path $k \in \{1, 2, 3\}$ and $(\alpha \cdot \eta)_{ik}$ is the interaction. A significant interaction validates that closure paths are not merely substitutable measurements of the same underlying quantity. Because \texttt{scipy.stats}--\texttt{statsmodels} version drift blocked the standard OLS route on the analysis host, we implement Eq.~\ref{eq:anova_interact} directly using QR-based least squares and Type-I sequential sums of squares in pure NumPy.

\subsubsection{Pairwise cell comparisons}

Pairwise mean differences between cells $i$ and $k$ are tested via Tukey's honest significant difference:

\begin{equation}
  q_{ik}
  \;=\;
  \frac{\bar{Y}_i - \bar{Y}\emph{k}{\sqrt{MS}{\textsf{within}} / n_h}},
  \label{eq:tukey}
\end{equation}

where $n_h$ is the harmonic mean of per-cell sample sizes (mono cells have $n = 150$; hetero and null cells have $n = 75$). A pair is significant at family-wise $\alpha = 0.05$ after Bonferroni correction over the $20 \times 19 / 2 = 190$ pairwise comparisons.

\subsubsection{Judge--metric correlation}

The Pearson correlation between the automated metric $Y$ and the judge rating $J$ per path,

\begin{equation}
  r_p \;=\; \frac{\textstyle\sum_{r} (Y_r - \bar{Y})(J_r - \bar{J})}{\sqrt{\textstyle\sum_r (Y_r - \bar{Y})^2 \cdot \textstyle\sum_r (J_r - \bar{J})^2}},
  \label{eq:corr}
\end{equation}

quantifies the extent to which each path's automated metric agrees with the judge SLM's semantic-equivalence rating. Rows with parse failures ($J = -1$) are excluded.

\subsubsection{Mixed-effects logit for per-stage decomposition}

To isolate which stage assignment drives closure, we treat each cell as the product of three stage choices and fit a generalised linear mixed model. Let $\text{closure}_{ij} \in \{0, 1\}$ denote whether function $j$ closed under cell $i$. We model

\begin{equation}
  \text{logit}\bigl(\Pr[\text{closure}_{ij} = 1]\bigr)
  \;=\;
  \beta_0
  + \beta_1 \cdot a_i(\textsf{spec})
  + \beta_2 \cdot a_i(\textsf{test})
  + \beta_3 \cdot a_i(\textsf{code})
  + \gamma_j,
  \label{eq:mixedlogit}
\end{equation}

where $\gamma_j \sim \mathcal{N}(0, \sigma_{\textsf{sample}}^2)$ is a random intercept for sample. Estimates $\hat{\beta}_1, \hat{\beta}_2, \hat{\beta}_3$ quantify each stage's contribution to closure odds.

\subsubsection{Judge--human agreement (reported in Section~\ref{sec:human_eval_addendum})}

For each binarised rating (judge threshold $J \geq \rho$ versus $J < \rho$), Cohen's $\kappa$ is

\begin{equation}
  \kappa
  \;=\;
  \frac{p_o - p_e}{1 - p_e},
  \label{eq:kappa}
\end{equation}

where $p_o$ is the observed agreement proportion between the judge SLM and a human annotator and $p_e$ is the chance-expected proportion under independence. Values follow Landis and Koch's bands: $\kappa > 0.6$ substantial, $\kappa > 0.8$ almost perfect. For the three-annotator (judge + two human) setting, we additionally compute Krippendorff's $\alpha$ for ordinal ratings $r_{ij}$ from rater $i$ on item $j$:

\begin{equation}
  \alpha
  \;=\;
  1 \;-\; \frac{D_o}{D_e},
  \quad
  D_o \;=\; \tfrac{1}{n}\sum_{j} \tfrac{1}{\binom{m_j}{2}}
            \sum_{i_1 < i_2} \delta(r_{i_1 j}, r_{i_2 j})^2,
  \label{eq:alpha}
\end{equation}

where $\delta$ is the ordinal-difference metric and $D_e$ is the expected disagreement under random pairing. Values $\alpha > 0.667$ are considered acceptable for inferential use. Both $\kappa$ and $\alpha$ are reported in the human-evaluation addendum (Section~\ref{sec:human_eval_addendum}); the headline numbers on the 60-pair three-annotator sample are $\alpha_{\textsf{ord}} = 0.376$ across the three human annotators and $\kappa_{\textsf{quad}} = 0.138$ between the judge SLM and the majority human rating.

\subsubsection{Sensitivity analysis}

To defend against threshold-choice attacks, closure rates are computed and reported at multiple $(\tau, \rho)$ combinations. For the mutation kill rate on Path~1, $\tau \in \{0, 0.5, 0.8\}$; for the judge threshold, $\rho \in \{2, 3, 4\}$. Ordinal stability of the per-cell ranking across thresholds is quantified using Kendall's $\tau_{\textsf{rank}}$ between the reference threshold $(\tau=0, \rho=3)$ and each alternative.

\subsubsection{Post-processing categorisation}

Every closure row is post-hoc labelled with the Algorithm~\ref{alg:validity} decision reason and the Algorithm~\ref{alg:test_filter} filter reason. The distribution of decision reasons per path (Section~5.2) is our primary judge--metric disagreement diagnostic; the distribution of filter reasons per (cell, path) triple (Section~5.3) is our primary diagnostic for the phi-4-in-test-stage attribution.

\subsubsection{Capability preservation}

For each path, we identify the strongest mono baseline (M4 for Paths~1 and~3; M3 for Path~2) and restrict analysis to samples on which that baseline achieves valid closure. For every other cell, the \emph{capability-preservation rate} is the fraction of those same samples on which the cell also achieves valid closure. Values near~$1$ indicate that the cell does not sacrifice mono strengths; substantially lower values expose cells where specialisation costs capability on samples the strongest mono handles cleanly.


\emph{End of Section~3 Methods.} Section~4 presents results anchored to Eqs.~\ref{eq:killrate}--\ref{eq:corr} and Table~\ref{tab:doe_summary}, organised by research question.
